\section{演化临界点：工具理性的危机}

\subsection{替代的非同步性}

把介质分析和验证结构放在一起，可以得到一个关于就业的直接推论，而且这个推论和流行叙事恰好相反。物理层面的黏滞延缓了 AI 对实体劳动的替代进程。搬砖、装配、维修、护理、驾驶、巡检这些工作，处在原子域，受质量、能量、摩擦和现场不可预测性保护，短期内难以被低成本替代。逻辑层面的低黏滞则加速了 AI 对认知劳动的替代。写作、翻译、整理、初级编程、常规分析、标准文档处理，处在比特域，恰好是生成式 AI 最擅长的区域。这种非同步性造成了一种社会性分离：\key{肉身还留在旧世界，职业尊严和智力上的自我评价已经被提前改写。}

\subhead{物理防线：能效剪刀差与容错性}

这道防线首先是成本曲线。部署机器人要支付采购、集成、折旧、维修、电力、监督与停机成本。雇佣人也有工资、福利、培训和安全支出；人的成长与维持成本，有些由家庭、社会和本人承担，没有进入雇主账本。低工资可能来自成本转移。一个岗位暂时保留下来，说明此刻的部署比较还没有翻转，不能据此认定人天然便宜。

最后一米常常聚集最不规则的工作。包裹路线可以排好，入口却可能找错，门禁可能失灵，收件人可能不在；仓库里的容器会变形，物体会偏离预定姿态。自动化正在吞下越来越多这样的情况，每增加一种，都要让工程可靠性和运营账本一起成立。用一幅永远失败的机器人漫画解释今天，和用一次演示宣布全面替代，是同一种偷懒。

最新进展恰好把两种时钟放在同一家公司里。Waymo在2026年10月27日宣布，接下来一周将在五个美国市场合计扩展近250平方英里的服务区域；10月20日公布的马德里计划仍从人工驾驶和当地技术验证开始，目标是2027年迎接首批乘客；10月16日公布的多伦多计划，则把首批乘客目标放在2028年。\cite{waymo2026expansion,waymo2026madrid,waymo2026toronto} 已有运营的扩区、当地验证与未来上线目标，不能折叠成同一项能力指标。图\ref{fig:deployment-stages}把这几种状态放在一起。

\begin{figure}[ht]
\centering
\begin{tikzpicture}[x=1cm,y=1cm,font=\small,
box/.style={draw=slate!45,rounded corners=2pt,align=left,inner sep=6pt,text width=3.45cm,minimum height=1.65cm},
arr/.style={-{Stealth[length=2mm]},line width=0.8pt,slate}]
\node[box,fill=deepblue!8] (a) at (0,0) {\textbf{能力}\\任务、场景与性能\\模型、车辆与工具};
\node[box,fill=terra!9] (b) at (4.6,0) {\textbf{当地部署}\\道路验证、设施、权限\\维护与应急衔接};
\node[box,fill=forest!9] (c) at (9.2,0) {\textbf{持续运营}\\实际服务与扩区\\成本、故障与后果};
\draw[arr] (a)--(b);
\draw[arr] (b)--(c);
\draw[arr] (c.south) -- ++(0,-0.5) -| (b.south);
\node[align=center,font=\scriptsize,text=slate] at (4.6,-1.95) {马德里：2026-10-20公布筹备，目标2027年迎客\\多伦多：2026-10-16公布筹备，目标2028年迎客};
\node[align=center,font=\scriptsize,text=slate] at (9.0,-2.95) {美国五个市场：2026-10-27公布\\未来一周合计扩区近250平方英里};
\end{tikzpicture}
\caption{能力、部署与运营，是三种不同的进展。Waymo案例来自2026年10月公告；扩区与上线年份分别按当时公布的计划记录，箭头表示需要衔接的工作。}
\label{fig:deployment-stages}
\end{figure}


把两边的成本写成两条曲线，替代发生的条件就是这两条曲线的交点，

\begin{equation}
C_{\text{人}}=w,
\qquad
C_{\text{机}}=\frac{K}{L}+m
\end{equation}

式中$w$为分配到人力方案上的年化成本，$K$为机器采购与安装投入，$L$为折旧年限，$m$为年度运行维护支出。在产出、质量、可用性与风险相当的前提下，较低的年化成本会推动替代。融资、利用率、切换支出与责任分配还要另行计入。账上更划算之后，采购、安装、认证和人员转岗仍要一项项发生。

这道防线有明确的边界。剪刀差保护的是劳动的相对价格，不是劳动本身。一旦电力成本出现数量级下降，比如可控核聚变真正落地，100 瓦对 3000 瓦的比较就会改变；一旦机器人的产量上来，资本支出也会被摊薄；一旦某个场景被充分结构化，容错性的溢价同样会消失。保护人类岗位的是两套成本曲线的相对位置，而不是某条不可逾越的物理定律。这一点必须说清楚，否则物理保护会被误读成永久豁免。

\subhead{心理崩塌：从万物之灵到廉价接口}

物理上的安全没有带来安心，反而制造了新的问题。这解释了一种普遍的困惑。许多人没有失业，却已经开始怀疑自己工作的意义。他们的岗位还在，是因为所在系统的高黏滞保护了它；但他们的技能被定价的方式，已经被低黏滞领域的生成能力改写。危险不是立刻发生的替代，而是价值证明接口的提前失效。

这里有一个容易混淆的地方。非同步性并不意味着认知劳动会先消失，而是意味着它的证明方式先失效。工作还在，定价方式变了；岗位还在，但进入这个岗位所需要的训练不再自动换回位置。把这两件事分开，才能理解为什么失业率这类数据往往看不出这场变化，而从业者的感受却很剧烈。数据统计的是岗位的有无，感受衡量的是位置是否还被承认。

要理解这种失效，需要回到旧的价值公式。自启蒙运动和工业革命以来，现代文明的价值观建立在一个等式上：人的价值等于人的社会功能，等于人掌握的技能。教育强调学好数理化，职场鼓励提升专业度，社会地位按产出效率分配。这套排序默认，越需要复杂脑力的工作越高等，越纯粹的体力工作越低等。大模型对这一排序的冲击是釜底抽薪。逻辑推理、编程、写作、绘画、法律咨询，这些曾被视为智慧领域的能力，变成了可以按量购买的廉价商品。画师磨练十年的光影技巧，生成模型一秒就能给出；程序员积累五年的架构经验，补全工具瞬间就能产出。

技能贬值的后果与收入无关。这种贬值带来的恐惧不是“我没饭吃了”，因为送外卖仍然可以糊口，而是“我没用了”。对中产和知识工作者而言，智力自尊被剥夺比收入下降更难承受。过去，一个人会写、会算、会总结，本身就说明他受过训练、接触过对象、在共同体中有位置；现在，这些能力被外部系统复制、放大和自动化，能力本身不再自动证明价值。技术外化的不是某一项孤立能力，而是旧社会用来识别训练、经验、资格和责任的接口；接口一旦被接管，人就必须迁移到别的接口上去证明自己。

在物理不失业和心理失能的夹缝里，一种新的工作形态已经成形。既然智能有清晰的规划能力却没有灵活的手，最高效的系统架构就是把规划交给算法，把动作交给人。外卖骑手和网约车司机是这个架构的早期样本：算法决定路线、定价、时间和奖惩，人负责拧动油门、爬楼梯、把餐送到门口、承受客户的情绪。人处理的是物理世界最后一段摩擦力。这种工作能带来收入，甚至不算低，但它不提供主体性。人不再是通过劳动创造价值的主体，而是被算法调用的一个部件，暂时比机械臂便宜；相比之下，工业时代的工人还是操作机器的主人。

Yuval Harari 提出过“无用阶级”的概念。在这个框架里它需要一处修正：问题不是经济上无用，事实上人仍然很实用，而是叙事上无用。文明的核心叙事一直是人类主导的探索、创造、征服；当最前沿的科学发现可以由系统完成，最复杂的调度由系统决策，人就成了旁观者。如果我看过的书系统都看过而且记得更牢，如果我写出的诗系统都能模仿而且韵律更准，如果我做出的决策系统都能优化而且更稳，那么我是谁，我存在的必要性在哪里。这种恐惧不像饥饿那样剧烈，但它持续，并且穿透力强。而它的现实后果，往往不是失业，而是一整代人撤离了攀登技能树这件事。

这一节给出的图景有两个面。物理正面，熵增定律和高黏滞把人暂时安全地留在搬运原子、对抗摩擦、伺候肉身的高熵劳动里。心理背面，比特世界的光速迭代抽掉了对自身位置的旧解释。两句话可以概括：不会失业，但会失魂。肉体被固定在旧世界的现场，精神却漂泊在被算法接管的地方。这比单纯的生计问题更难处理，因为它没有一个可以直接对应的补偿方案，也因为它发生在一个更大的系统变化里：效率的提升本身正在制造新的稀缺。

\subsection{效率通货膨胀与发展的结构性瓶颈}

效率提升释放的空间，很快可能被新的需求填满。生成变便宜，方案、代码和草稿一起涌来；审核与维护若跟不上，人就从起草者变成排队处理候选的审核员。工作并非依守恒定律注定回来，而是到达速率超过了接收能力。真正要追的是：省下来的那一步，给下一步增加了多少东西。


\begin{figure}[ht]
\centering
\begin{tikzpicture}[x=0.7cm,y=0.50cm,font=\scriptsize]
\fill[terra!8] (3,0) rectangle (7,5.8);
\draw[-{Stealth[length=1.6mm]}] (0,0)--(15.8,0) node[right]{时间};
\draw[-{Stealth[length=1.6mm]}] (0,0)--(0,5.8) node[above right]{入队与处理速率};
\draw[deepblue,line width=1.2pt] (0,1)--(3,1)--(3,4)--(15,4);
\draw[terra,dashed,line width=1.2pt] (0,2)--(7,2)--(7,5)--(15,5);
\node[deepblue,anchor=south west] at (8,4.05){入队速率 $\lambda$};
\node[terra,anchor=south west] at (8,5.05){处理能力 $\mu$};
\node[align=center] at (1.5,3.3){尚有余量};
\node[align=center] at (5,3.1){入队超过处理\\积压开始增长};
\draw[slate!50,densely dotted] (3,0)--(3,-1.1);
\draw[slate!50,densely dotted] (7,0)--(7,-1.1);
\node[anchor=north,align=center] at (3,-0.15){生成提速\\$t_g$};
\node[anchor=north,align=center] at (7,-0.15){验证扩容\\$t_v$};
\begin{scope}[yshift=-5.5cm]
\fill[forest!16] (3,0)--(7,4)--(15,0)--cycle;
\draw[-{Stealth[length=1.6mm]}] (0,0)--(15.8,0) node[right]{时间};
\draw[-{Stealth[length=1.6mm]}] (0,0)--(0,4.9) node[above right]{待处理候选 $Q$};
\draw[forest,line width=1.2pt] (0,0)--(3,0)--(7,4)--(15,0);
\draw[slate!55,densely dotted] (7,0)--(7,4);
\node[anchor=south,align=center] at (7,4.15){扩容时积压最多\\能力补上，队列仍要消化};
\node[anchor=north] at (3,-0.15){$t_g$};
\node[anchor=north] at (7,-0.15){$t_v$};
\node[anchor=north,align=center] at (15,-0.15){$t_c$\\积压清空};
\node[align=center,text=slate] at (8,-1.55){$Q>0$时：$\dot Q=\lambda-\mu$\\斜率由两种速率之差决定，清空时间由积压量决定};
\end{scope}
\end{tikzpicture}
\caption{瓶颈怎样移动。图中$\lambda$是进入指定验证队列的候选速率，$\mu$是同口径的处理能力。生成提速后$\lambda>\mu$，候选积压；验证扩容后$\mu>\lambda$，队列逐步消化。曲线为示意，纵轴不表示实测值。排队数量服从流量收支，社会劳动总量不由这个式子决定。}
\label{fig:bottleneck-migration}
\end{figure}


技术进步常常从一个局部降阻开始。一条路变短，一个设备变高效，一个模型变便宜，一个流程被自动化，一个指标变清楚，一个风险被安全装置压低。直觉上，局部阻力下降，整体系统应该变好。成本降低，效率提高，选择增多，风险减少，控制增强。但复杂系统经常给出相反结果。多一条路，交通可能更堵；效率提高，总消耗可能上升；安全增强，行为可能更冒险；指标更明确，真实目标可能被腐蚀；模型越被相信，被预测对象越可能改变；自动化减少常规任务，却把人推向更稀有、更困难、更高责任的异常场景。这些现象有不同名字。

\subhead{局部降阻与主体响应}

这些现象有不同的名字，分布在能源、交通、管理、预测、安全和自动化六个互不相干的领域。把它们并排放在一起，共享的结构会自己显示出来。

\begin{table}[ht]
\centering
\caption{六类悖论共享的结构：局部降阻、主体响应、全局重排}
\label{tab:ch1-paradox}
{\small
\begin{tabularx}{\textwidth}{@{}l >{\RaggedRight\arraybackslash}p{2.7cm} >{\RaggedRight\arraybackslash}p{3.4cm} >{\RaggedRight\arraybackslash}X@{}}
\toprule
悖论 & 局部降阻 & 主体响应 & 全局结果 \\
\midrule
杰文斯 & 单位使用成本下降 & 更多场景接入，使用次数上升 & 总消耗不降反升 \\
布雷斯 & 新增一条可用路径 & 每个出行者按局部最优改道 & 均衡落点变差，全网通行时间上升 \\
古德哈特与坎贝尔 & 指标变得清晰可测 & 被测方围绕指标重组行为 & 指标脱离原目标，测量失效 \\
卢卡斯与反身性 & 模型给出更准的预测 & 被预测方依据模型调整自身 & 原关系失效，系统轨道改变 \\
佩尔兹曼 & 单次失效概率下降 & 主体提高速度、频率或杠杆 & 总风险暴露扩张，净收益被抵消 \\
自动化悖论 & 常规操作被自动接管 & 人退守到稀有的异常场景 & 剩余任务更难，责任更集中 \\
\bottomrule
\end{tabularx}}
\end{table}

表 \ref{tab:ch1-paradox} 里六行的第二列各不相同，第四列却落在同一类后果上，这不是巧合。每一行都可以读成同一个三步过程：某个局部条件被改变，处在系统里的主体按自己的利益重新选择，总量在新的均衡上重排。六行后面逐个处理；需要先记住的是第三列：技术改变的是主体的选择场，不是主体的行为规则，而前者一变，后者就必然跟着变。这也解释了为什么这类悖论总是出现在“局部指标改善”和“整体状况恶化”同时为真的报告里：两者测的不是同一层的量，把两层混在一起看，得到的只会是矛盾。

技术不是简单从系统中减去一块阻力。它更像改变了系统的成本地形。某个局部势垒被削低之后，主体不会静止不动，而会重新选择路径。主体重新选择之后，流量场、风险场、指标场、任务分布和责任分布都会重排。系统最后不是停在原状态上享受局部收益，而是在新的边界条件下形成新的均衡。悖论出现于这里：局部势垒下降了，但新均衡下的总黏滞更高。

所以，技术效果不能只看局部参数。一个局部更快、更省、更准、更安全，并不自动推出整体更快、更省、更准、更安全。复杂系统的全局代价不是局部收益的线性相加，而是主体响应之后，由流量、瓶颈、风险、验证和责任共同生成的结果。可以把这个统一框架概括为：局部降阻改变成本地形，主体响应重排系统分布，系统进入新均衡，隐性瓶颈在新位置展现。

沿着这个框架往下看，六类悖论可以分成三组。第一类，是杰文斯悖论：低阻膨胀。杰文斯悖论最初讨论能源效率。一个设备更省煤、更省油、更省电，直觉上应该减少总消耗。但当效率提高降低了单次使用成本，更多人会使用，更多场景会接入，更多需求会被释放，最后总消耗可能不降反升。\cite{jevons1865coal,allwood2012sustainable}

这不是能源史中的特殊反直觉，而是通量反弹。一个通道的阻力下降后，系统中会有更多流量被吸引过来。原来因为贵而不做的事，现在会被大量做；原来因为慢而慎用的能力，现在会变成常规调用；原来因为稀缺而受到约束的资源，现在会被嵌入更多场景。

效率降低的是单位阻力，放大的可能是总通量。AI 时代，杰文斯悖论更加明显。生成更便宜，于是内容更多；代码更容易生产，于是软件复杂度和维护负担增加；方案更容易生成，于是验证压力上升；算力更容易调用，于是总算力需求继续膨胀。技术降低了单次操作的阻力，却让更多操作穿过系统。因此，效率不是总成本的自动反面。一个系统是否真正减负，不取决于单次操作是否更便宜，而取决于总通量扩大之后，能源、散热、验证、维护、责任和吸收能力是否仍然承受得住。杰文斯悖论在技术黏滞中的位置，就是说明：局部黏滞下降，可能诱发全局通量上升。把单位成本与使用总量之间的关系写成弹性形式，

\begin{equation}
E=q\,N(p),\qquad p=\alpha q,
\qquad
\varepsilon\equiv-\frac{d\ln N}{d\ln p},
\qquad
\frac{d\ln E}{d\ln q}=1-\varepsilon
\end{equation}

式中$q$为单次操作的物理消耗，$p$为单价，$N(p)$为单位时间的操作次数；这里用固定$α$把两者联系，$E$是单位时间总消耗。弹性$ε<1$时，效率改进降低总消耗；$ε>1$时，新增需求超过单次节省。2026年4月，国际能源署同时记录了单项AI任务能效的显著改善和2025年数据中心总用电17\%的增长。\cite{iea2026questions} 单次更省与总量上升可以并存；具体反弹幅度仍须由需求数据识别。

第二类，是布雷斯悖论：新路径坏均衡。布雷斯悖论说明，在某些交通网络中，多修一条路反而会让所有人的通行时间变长。问题不在于那条路本身有害，而在于它改变了每个司机的选择场。每个人按局部最优重新选择路线，流量重新分布，最后系统落入更差均衡。\cite{braess1968}

这不是交通系统的特殊现象，而是拓扑耦合的普遍问题。路径成本不是路径本身的静态属性，而是流量加载之后生成的属性。一条路在少数人使用时可能低阻，在所有人涌入后就可能成为新的拥堵点。把每个出行者的最小化和整个系统的最小化并排写出，

\begin{equation}
f_i>0\ \Longrightarrow\ c_i(\mathbf f)=\min_{j\in\mathcal P}c_j(\mathbf f);
\qquad
\min_{\mathbf f}\sum_{e\in E}f_{e}\,c_{e}(f_{e})
\end{equation}

式中$f_e$为路段流量，$c_e(f_e)$为对应通行时间，$\mathcal P$为同一起终点的路线集合，路线成本由沿途路段相加。固定需求下，用户均衡要求每条被使用的路线都达到当前最小成本；系统最优则最小化全网总通行时间。新增道路扩大可行集合，允许不用新路时，系统最优不会变差；独立选择形成的新均衡却可能更差。布雷斯悖论发生在这两者的缝隙里：选项增加了，协调并没有随之完成。

一个工具在少数人使用时是优势，被集体复制后却可能成为新的瓶颈。新增道路扩大物理上的可达集合，使用者共同选择之后形成什么流量，则是另一回事。技术提供路径，主体选择路径，压力在选择之后重新分布。软件和组织系统也值得追问这个问题，具体结果由各自机制决定。

\subhead{测量与模型}

通量与路径的重排发生在物质流动层面，而这些重排要被主体感知、比较和选择，还必须经过信息层，信息层本身也会被这一轮重排改写。第三类，是古德哈特定律与坎贝尔定律：指标腐蚀。古德哈特定律通常被概括为：当一个指标成为目标，它就不再是好指标。\cite{goodhart1975problems} 坎贝尔定律也指出，社会指标越被用于决策，越容易被腐蚀，并扭曲它本来要监测的过程。\cite{campbell1979assessing}

这类悖论的根源，是测量进入了系统。指标本来像观测量，用来近似描述某个复杂真实。但当它进入奖惩、资源分配和组织控制之后，它就不再只是观测量，而变成控制量。被测对象会学习指标、迎合指标、套利指标，甚至牺牲真实目标来优化指标。于是，指标与目标之间的关系被行为反馈改写。学校围绕排名调整教学，医院围绕绩效调整服务，平台创作者围绕推荐算法调整内容，企业围绕 KPI 调整行动。指标越重要，行为越围绕指标重组；行为越围绕指标重组，指标越可能脱离原本要衡量的真实目标。

从物理机制上看，指标不是无扰动测量。它一旦接入激励结构，就改变了系统的自由度分布。它把复杂真实压缩成少数可优化变量，同时诱导主体把行动集中到这些变量上。指标越强，真实越可能被指标重新塑形。

因此，验证系统本身也会制造黏滞。真正成熟的验证，不只是提高测量精度，还要防止测量通道被行为反向占用。否则，指标越清楚，系统越可能围绕指标形成新的低阻路径，而真实目标反而被屏蔽。古德哈特定律和坎贝尔定律出现在测量这一侧：局部测量清晰，可能诱发全局目标腐蚀。

第四类，是卢卡斯批判与反身性：模型进入系统后改变对象。卢卡斯批判指出，用过去行为关系建立的模型，一旦被用于政策，主体会调整行为，原来的关系就可能失效。\cite{lucas1976econometric} 反身性理论则强调，在市场、组织和社会系统中，参与者的认知和行动会改变被认知对象本身。\cite{soros1987alchemy}

这类问题的根源，是模型进入行动回路。模型本来用来描述系统，但当模型被用来决策，系统中的主体会根据模型改变行为。市场预测会影响交易者行为，风险模型会影响金融机构策略，政策模型会影响企业和公众预期，组织诊断会影响组织成员反应。模型一旦进入系统，就不再只是外部表征，而成为系统内部变量。

这和物理对象不同。杯子不会因为你建模它而主动改变自身；但市场、组织、舆论和谈判对手会。它们会理解、适应、反制和利用模型。模型越被相信，越可能改变它预测的对象；预测越广泛进入行动，系统越不再沿原来的轨道演化。这类系统的困难，不在于模型不够准确，而在于模型输出会改变边界条件。原来模型面对的是某个状态下的系统；模型被使用后，系统已经不再是原来的系统。

预测进入行动回路之后，会成为新的扰动源。它可能引发自我实现，也可能引发自我挫败。所以，高耦合系统中的模型不能只追求静态准确。它还必须追问：模型被使用之后，会怎样改变系统自身？模型不是只在系统外画图，它会进入流场，改变下一轮流向。卢卡斯批判和反身性把问题推进到预测层面：局部预测能力提升，可能诱发被预测对象的重构。

\subhead{风险与任务分布}

测量与模型被改写之后，主体对世界的判断会变，判断变了，行动方式本身也会变。第五类，是佩尔兹曼效应：风险补偿。佩尔兹曼效应说明，安全技术降低局部风险之后，主体可能采取更冒险行为，抵消一部分安全收益。道路更安全，驾驶者可能更快；辅助驾驶更强，驾驶者可能更分心；金融风控更细，机构可能更激进；平台审核更自动化，用户可能更频繁地测试边界。\cite{peltzman1975effects}

这不是说安全无用，而是说风险不是静态物质。风险是技术、主体和环境共同生成的分布。安全装置降低了某个局部失效概率，但主体会根据新的安全边界重新调整暴露量、频率和冒险强度。单次风险下降，总风险未必同比下降。

这里的关键是边界外扩。系统获得更多安全余量后，主体可能把这部分余量转化为更高速度、更高频率、更高杠杆、更低注意力或更激进行为。安全降低了局部势垒，却也改变了行动空间。最终总风险是否下降，取决于局部风险压低的幅度，是否大于行为边界外扩带来的新增暴露。所以，风险治理不能只计算单次失效概率，还要计算主体响应后的总暴露。否则，安全系统会高估自己的净收益。降低风险不等于消灭风险；它可能改变风险被使用的方式。佩尔兹曼效应是这套逻辑在风险领域的版本：局部安全性提高，可能诱发全局风险暴露扩张。把这次比较写成式子是直接的，

\begin{equation}
R=p\,n\,d,
\qquad
\Delta R\approx n\,d\,\Delta p+p\,d\,\Delta n,
\qquad
\Delta R<0\ \Longleftrightarrow\ \frac{|\Delta p|}{p}>\frac{\Delta n}{n}
\end{equation}

$R$ 为单位时间内的期望损失，单位是损失每单位时间，例如事故数每年；$p$ 为单次暴露发生失效的概率，无量纲；$n$ 为单位时间内的暴露次数，单位是次每年；$d$ 为单次失效的平均损失，单位是损失每次。安全装置使 $p$ 下降（$\Delta p<0$），主体响应使暴露量上升（$\Delta n>0$），一阶近似下两者对 $R$ 的作用方向相反，总风险下降的条件是相对减小的幅度大于相对增加的幅度。正文所说的“局部风险压低的幅度”对应 $|\Delta p|/p$，“行为边界外扩带来的新增暴露”对应 $\Delta n/n$：压低的幅度必须大于外扩的幅度，安全才有净收益，这也是风险治理不能只计算单次失效概率的原因。它的适用边界是：一阶近似只在两个相对变化都较小时成立，当主体响应是离散的（干脆关闭安全装置，或者退出这类活动），或者单次损失 $d$ 本身随暴露强度变化时，这个线性分解就失效。这个式子还只回答总风险的升降，不回答风险在人群中的分布。安全装置常常只对其中一部分人有效，外扩后的暴露由另一部分人承担。

第六类，是自动化悖论：残差集中。自动化悖论，也常被称为自动化的反讽，指的是：自动化减少了常规操作，却可能让人类在异常时更难接管。机器处理常规流程，人类远离过程细节；一旦异常出现，人类既缺少连续参与，又必须快速判断。常规负担下降，剩余负担却更高密度、更难形式化、更难追责。\cite{bainbridge1983ironies}

这不是负担消失，而是任务分布改变。自动化通常先吸收高频、规则化、可重复的部分，把低频、边缘、模糊、责任重的部分留给人。于是，人类不再承担连续操作负担，却承担尾部残差负担。任务数量减少，任务密度上升；日常负担下降，异常责任上升。系统看起来更自动化，人的角色却被推向更难的位置。

AI 越强，这个问题越突出。模型可以处理大量常规判断，但剩下的问题往往更难验证、更难解释、更难负责。自动化降低了前台操作黏滞，却把残差集中到异常、验证和责任环节。它不是把人从系统中完全释放出来，而是重新安排人类在系统中的位置。自动化悖论在技术黏滞中的位置，就是说明：局部任务自动化，可能诱发残差负担向尾部集中。

\subhead{阻力场的重分布}

这些悖论名称不同，场景不同，但底层结构相同。杰文斯悖论是通量响应。布雷斯悖论是路径响应。古德哈特定律和坎贝尔定律是指标响应。卢卡斯批判和反身性是模型响应。佩尔兹曼效应是风险响应。自动化悖论是任务分布响应。

它们共同构成一种技术悖论范式：技术先降低局部势垒；主体根据新的势垒重新行动；行动改变流量场、风险场、测量场、任务分布和责任分布；系统在新的边界条件下形成新均衡。悖论出现于：新均衡中的总代价，超过了局部降阻带来的收益。这不是创造新的悖论名称，而是把已有悖论放进同一个机制框架。

因此，判断一项技术是否真正降低系统黏滞，不能只看局部效率、局部成本、局部速度、局部安全或局部准确率。必须看它进入系统后，是否改变了总通量、均衡状态、指标行为、模型反身性、风险暴露、验证压力和责任结构。技术悖论不是反技术。它只是提醒我们：\key{局部优化不是系统优化。}开放系统不会被动接受收益，它会响应收益、放大收益、套利收益，并把收益重新转化为新的压力。一个技术越有效，越可能诱发行为重分布；一个低阻路径越好，越可能吸引更多流量；一个指标越有用，越可能被围绕它优化；一个模型越被相信，它预测的对象就越会因被预测而改变。技术黏滞的深层意义也在这里。

黏滞不是一堵固定的墙，而是一种会随技术输入重新分布的阻力场。技术降低一个位置的阻力，系统可能在另一个位置展现新的阻力。阻力可以从能源转移到散热，从生成转移到验证，从路径转移到均衡，从安全转移到行为，从自动化转移到残差，从模型转移到责任。

真正成熟的技术判断，不是问“这个局部是否更快、更省、更准、更安全”，而是问：这个局部降阻之后，成本地形如何改变？主体会怎样重新行动？流量会推向哪里？瓶颈会在哪里重新展现？风险和责任会由谁承担？系统会进入什么新均衡？总黏滞到底下降了，还是只是换了位置？

能源与基础设施还在给这场重排设定条件。获取能源需要投入能源，净剩余取决于技术路线、地点与核算边界；材料会老化，维护支出取决于资产状态和服务要求。不能用一条笼统的历史下降曲线替所有能源技术判定未来。算力许诺的扩张，仍要落到具体的电源、设备和维护能力上。

技术悖论的共同形式，正是局部降阻之后的全局重结算。局部降阻不是终点。新均衡才是答案。

但这一步目前还只是定性的。六类悖论各有自己的术语、场景和文献传统，它们能被放进同一个框架，是因为底层结构相同，而不是因为它们已经由同一组变量描述。局部降阻之后，代价落在通量、均衡、测量、风险还是任务分布上，在现有表述里仍然只能逐案回答。如果这个共同形式是成立的，那么阻力的转移、通量的响应和新均衡的形成之间，就应当存在一组可以被重复使用的量，而不只是名词上的归类。这一点还没有被写出来。

\subsection{终极追问：后工具时代的悖论}

上一节把阻力场的重分布讲成了一件纯结构的事，但阻力落到人身上从来不平均。降低的那部分阻力从某些岗位上被抽走，新增的那部分阻力被压到另一些岗位上。谁承担转移后的黏滞，是分配问题，不是效率问题；而在讨论分配之前，需要先把这类判断的边界说清，否则技术黏滞论很容易被读成另一种情绪。

\subhead{技术效果的总账}

数字支付、导航、工业仿真、自动化实验和代码工具，都在不同环节降低了转移成本。评价这些改进，需要同时考察输出、执行、验证、维护与失败处理。哪些成本减少，哪些负担转移，哪些风险尚未暴露，应按系统边界分别记录。

它也不是万能解释。现实技术落地还受到政治、文化、资本周期、商业模式、组织能力、利益结构和社会心理影响。技术黏滞不是要把一切社会问题还原为资源、时间和拓扑，而是提供一个底层视角：当智能试图从表示空间进入现实系统时，它必然遇到介质阻力。这个阻力不是偶然噪声，而是现实系统本身的结构属性。AI 时代最重要的能力，因而不只是更强的生成能力，而是把生成结果转化为可执行、可验证、可维护、可担责现实改变的能力。未来的竞争也不只是模型能力之争，而是验证体系、工程体系、组织体系、基础设施和责任制度的竞争。谁能更好地管理黏滞，谁就能更有效地把智能转化为现实生产力。

\subhead{三重悖论}

边界说清之后，可以把本章的推演收成三个互相咬合的悖论。第一个是效率悖论。人曾经通过变得更快、更准、更强来证明价值，而这个方向的尽头是工具人自身的消亡：人越是在逻辑、计算和执行层面优化自己，就越逼近碳基大脑的生物极限，同时硅基系统在同一条赛道上以指数速度前进。在这条赛道上，终点线就是淘汰线。这不是能力的差距，是赛道的选择。

第二个是逃逸悖论。人试图通过虚拟化逃离物质世界的重量，结果是虚拟世界的自由度越高，对物理世界的依赖越深。越渴望零摩擦的体验，就越需要在原子世界里建设高能耗、高熵增的基础设施来支撑它。试图切断对物理现实的依赖，结果只是让依赖变得更硬。想飞升，反而加剧了坠落。

第三个是意义悖论，它是前两个叠加的结果。现代人的意义感建立在被需要之上，也就是工具价值。当效率悖论宣布人作为智能工具不再被需要，而逃逸悖论宣布无处可逃，意义的空位就出现了。人既失去了作为生产者的骄傲，又无法像古代贵族那样心安理得地只做体验者，因为物理基座仍然需要有人维护。人卡在有用的机器和自由的主体之间，两边都进不去。

\subhead{从怎么做转向是什么}

这三个悖论有一个共同前提：它们都假定问题出在能力上，只要跑得更快就能解开。如果这个前提不成立，那么技术层面的再技能化和经济层面的补偿都只是止痛，不是解法。当成为更好的工具这条路走到尽头，而逃往无重力的虚拟在物理上不可能，问题就必须换一个提法。原来的问题是：我们如何做得比机器更好。答案是不能。换过之后的问题是：当作的能力被外包给机器，是的内容是什么。

这里需要一次转向，而且这次转向不发生在技术层面。以效率为唯一坐标的那段历史，把人的位置定义在功能上；功能一旦可被廉价复制，位置就悬空了。新的位置不在功能里，而在别的地方：在能承受、能判断、能把后果接住这些事情上。这些能力不像算力那样可以集中采购，它们只能在具体的人身上、通过具体的时间积累出来。人类文明以模仿机器、打磨技能、追求效率来换取生存权的那段时间已经结束，接下来的问题不是在这张旧试卷上考得更高，而是重新定义考题。

这里有一个时间上的限定必须讲清楚。物理世界的滞后不是因为人不够努力，而是一段窗口期。在这段窗口里，人既要接受在效率上被机器超越，又要守住不被机器定义的底线。窗口不会一直开着，因为前面几节讲的三本账都在移动：能源成本可能下降，机器人的资本支出会被摊薄，高熵场景会被逐块结构化。缓冲的价值在于给了转向的时间，而不在于承诺了转向的结果。

转向的方向可以从历史里读出来。技术每外化一种能力，都会切断能力与价值之间的旧绑定。文字让记忆变轻，价值从背诵迁移到理解；印刷让复制变轻，价值从拥有文本迁移到解释和组织信任；蒸汽机让肌肉变轻，价值从做功迁移到控制、维护和调度；计算机让规则执行变轻，价值从执行迁移到问题定义；大模型让符号生成变轻，价值必须迁移到验证、准入和承受。每一次都是同一种结构：能力被外部系统接管，人退到更难伪造的位置。而这条线也说明，迁移从来不是自动发生的。文字出现之后，不是每个记忆者都自然成为解释者；印刷术出现之后，不是每个抄写者都自然进入出版制度；蒸汽机出现之后，也不是每个手工业者都自然变成工程师。迁移需要新的接口、新的制度，以及围绕这些接口的持续争取。

这条线还有一层限定，第六章会展开。迁移不是一条温和的上升曲线，而是一场博弈，高黏滞的位置并不天然属于人：今天被当作人的价值凭据的那些接口，同样可以被平台收编、被资本定价、被合同切割，也可以在专家系统和制度机器之间被重新分配。如果这些接口被少数机构垄断，人不会自动升格，只会从生产者变成被分类、被评分、被调参的对象。迁移也不可能只靠个人努力完成，它需要新的通道，让更多人能够表达真实经验、拒绝错误目标、参与后果的分配。

技术进步不会让世界变成无摩擦的信息空间。它会让某些环节变轻，同时让另一些环节变重；让某些过程加速，同时把瓶颈转移到验证、维护和责任；让表示空间迅速膨胀，同时迫使现实系统重新分配资源、时间和结构。理解 AI 时代，不能只看智能有多强，还要看智能如何穿过现实。真正持续有效的技术路径，属于那些能够识别黏滞、分配黏滞、降低不必要黏滞、保留必要黏滞，并在黏滞条件下完成可靠状态转移的人和制度。真正的未来，不属于无阻力的幻想。

但智能本身为什么会遇到这些阻力，还没有回答。要回答它，需要先问一个更早的问题：在一个由能量梯度驱动、由耗散维持的世界里，智能这种东西为什么会出现。这个问题不解决，前面所有的账目都只是一份清单。

\furtherreading{%
Kurt Gödel，\textit{Über formal unentscheidbare Sätze der Principia Mathematica und verwandter Systeme I}\\
Hans Moravec，\textit{Mind Children}\\
W. Stanley Jevons，\textit{The Coal Question}\\
Dietrich Braess，“Über ein Paradoxon aus der Verkehrsplanung”\\
Donald T. Campbell，“Assessing the Impact of Planned Social Change”\\
Sam Peltzman，“The Effects of Automobile Safety Regulation”\\
Lisanne Bainbridge，“Ironies of Automation”\\
Nassim Nicholas Taleb，\textit{Skin in the Game}}


% <<<<<<<<<< parts/book-ch01.tex <<<<<<<<<<


% >>>>>>>>>> parts/book-ch02a.tex >>>>>>>>>>
%% ===================== 本章参考文献 =====================
\printbibliography[heading=subbibliography,title=参考文献]
